% GENERATED FILE. Edit canonical lessons, then run npm run generate:books.
% canonical-source-hash: fnv1a64:81099c54a0858bf8
% canonical-lessons: KA-C175-beli, KA-C175-kamba, KA-C175-pakka, KA-C175-ace, KA-C175-ice, KA-R175-first-pass-words-from-chapters-167-170, KA-R175-second-pass-words-from-chapters-171-175

\chapter{Fence, Pillar, Beside, On the other side, On this side}
\label{ch:ka-fence-pillar-beside-on-the-other-side-on-this-side}
\hlchaptermodality{175}

\begin{chapteropening}
\textbf{By the end of this chapter:} \emph{I can say a fence, a pillar, beside, on the other side and on this side.}

Complete the last lesson of chapter 175: I can say a fence, a pillar, beside, on the other side and on this side.
\end{chapteropening}

\section[\texorpdfstring{\kn{ಬೇಲಿ} (bēli)}{ಬೇಲಿ (bēli)}]{\kn{ಬೇಲಿ} (bēli) — a fence}

\label{lesson:KA-C175-beli}



\begin{quote}
Before the new one: say the Kannada for a zoo, then the Kannada for a palace.
\end{quote}

\subsection*{You'll want to know: \kn{ಬೇಲಿ}}
\textbf{\kn{ಬೇಲಿ}} — \emph{bēli} — \textquotedblleft{}a fence\textquotedblright{}.

\begin{grammarlens}[title={What you've built}]
One more word for this chapter.
\end{grammarlens}

\subsection*{Your turn}
\begin{itemize}
\raggedright
  \item \emph{Say it:} \emph{bēli}
  \item \emph{Say it:} \emph{bēli}, once more
  \item \emph{Say it:} say \emph{bēli}
  \item \emph{Recall:} say the Kannada for a zoo, then the Kannada for a palace, then say \emph{bēli} again
\end{itemize}

\begin{tcolorbox}[breakable,colback=teal!4,colframe=teal!35!black,title={Before you move on}]
What does \kn{ಬೇಲಿ} mean? (\textquotedblleft{}A fence\textquotedblright{}.) Say it once more.
\end{tcolorbox}
\section[\texorpdfstring{\kn{ಕಂಬ} (kamba)}{ಕಂಬ (kamba)}]{\kn{ಕಂಬ} (kamba) — a pillar}

\label{lesson:KA-C175-kamba}



\begin{quote}
Before the new one: say the Kannada for a palace, then the Kannada for a fence.
\end{quote}

\subsection*{You'll want to know: \kn{ಕಂಬ}}
\textbf{\kn{ಕಂಬ}} — \emph{kamba} — \textquotedblleft{}a pillar\textquotedblright{}.

\begin{grammarlens}[title={What you've built}]
One more word for this chapter.
\end{grammarlens}

\subsection*{Your turn}
\begin{itemize}
\raggedright
  \item \emph{Say it:} \emph{kamba}
  \item \emph{Say it:} \emph{kamba}, once more
  \item \emph{Say it:} say \emph{kamba}
  \item \emph{Recall:} say the Kannada for a palace, then the Kannada for a fence, then say \emph{kamba} again
\end{itemize}

\begin{tcolorbox}[breakable,colback=teal!4,colframe=teal!35!black,title={Before you move on}]
What does \kn{ಕಂಬ} mean? (\textquotedblleft{}A pillar\textquotedblright{}.) Say it once more.
\end{tcolorbox}
\section[\texorpdfstring{\kn{ಪಕ್ಕ} (pakka)}{ಪಕ್ಕ (pakka)}]{\kn{ಪಕ್ಕ} (pakka) — beside, next to; the side}

\label{lesson:KA-C175-pakka}



\begin{quote}
Before the new one: say the Kannada for a fence, then the Kannada for a pillar.
\end{quote}

\subsection*{You'll want to know: \kn{ಪಕ್ಕ}}
\textbf{\kn{ಪಕ್ಕ}} — \emph{pakka} — \textquotedblleft{}beside, next to; the side\textquotedblright{}.

\begin{grammarlens}[title={What you've built}]
One more word for this chapter.
\end{grammarlens}

\subsection*{Your turn}
\begin{itemize}
\raggedright
  \item \emph{Say it:} \emph{pakka}
  \item \emph{Say it:} \emph{pakka}, once more
  \item \emph{Say it:} say \emph{pakka}
  \item \emph{Recall:} say the Kannada for a fence, then the Kannada for a pillar, then say \emph{pakka} again
\end{itemize}

\begin{tcolorbox}[breakable,colback=teal!4,colframe=teal!35!black,title={Before you move on}]
What does \kn{ಪಕ್ಕ} mean? (\textquotedblleft{}Beside, next to; the side\textquotedblright{}.) Say it once more.
\end{tcolorbox}
\section[\texorpdfstring{\kn{ಆಚೆ} (āce)}{ಆಚೆ (āce)}]{\kn{ಆಚೆ} (āce) — on the other side, over there}

\label{lesson:KA-C175-ace}



\begin{quote}
Before the new one: say the Kannada for a pillar, then the Kannada for beside.
\end{quote}

\subsection*{You'll want to know: \kn{ಆಚೆ}}
\textbf{\kn{ಆಚೆ}} — \emph{āce} — \textquotedblleft{}on the other side, over there\textquotedblright{}.

\begin{grammarlens}[title={What you've built}]
One more word for this chapter.
\end{grammarlens}

\subsection*{Your turn}
\begin{itemize}
\raggedright
  \item \emph{Say it:} \emph{āce}
  \item \emph{Say it:} \emph{āce}, once more
  \item \emph{Say it:} say \emph{āce}
  \item \emph{Recall:} say the Kannada for a pillar, then the Kannada for beside, then say \emph{āce} again
\end{itemize}

\begin{tcolorbox}[breakable,colback=teal!4,colframe=teal!35!black,title={Before you move on}]
What does \kn{ಆಚೆ} mean? (\textquotedblleft{}On the other side, over there\textquotedblright{}.) Say it once more.
\end{tcolorbox}
\section[\texorpdfstring{\kn{ಈಚೆ} (īce)}{ಈಚೆ (īce)}]{\kn{ಈಚೆ} (īce) — on this side}

\label{lesson:KA-C175-ice}



\begin{quote}
Before the new one: say the Kannada for beside, then the Kannada for on the other side.
\end{quote}

\subsection*{You'll want to know: \kn{ಈಚೆ}}
\textbf{\kn{ಈಚೆ}} — \emph{īce} — \textquotedblleft{}on this side\textquotedblright{}.

\begin{grammarlens}[title={What you've built}]
One more word for this chapter.
\end{grammarlens}

\subsection*{Your turn}
\begin{itemize}
\raggedright
  \item \emph{Say it:} \emph{īce}
  \item \emph{Say it:} \emph{īce}, once more
  \item \emph{Say it:} say \emph{īce}
  \item \emph{Recall:} say the Kannada for beside, then the Kannada for on the other side, then say \emph{īce} again
\end{itemize}

\begin{tcolorbox}[breakable,colback=teal!4,colframe=teal!35!black,title={Before you move on}]
What does \kn{ಈಚೆ} mean? (\textquotedblleft{}On this side\textquotedblright{}.) Say it once more.
\end{tcolorbox}
\section[(dialogue)]{Words from chapters 167-170 — a first pass}

\label{lesson:KA-R175-first-pass-words-from-chapters-167-170}



\begin{quote}
Say the words for on the other side and on this side.
\end{quote}

\subsection*{Your turn}
Say these aloud:

\begin{itemize}
\raggedright
  \item to drop off, to improve, to support, to burst, to snore
  \item the forehead, the throat, the heel, the thigh, the palm
  \item patience, jealousy, loneliness, curiosity, a feeling
  \item childbirth, eyesight, an infection, a blister, nutritious
\end{itemize}

\begin{tcolorbox}[breakable,colback=teal!4,colframe=teal!35!black,title={Before you move on}]
To drop off? (\textbf{\kn{ಉದುರು}}, \emph{uduru}.) To improve? (\textbf{\kn{ಸುಧಾರಿಸು}}, \emph{sudhārisu}.) To support? (\textbf{\kn{ಬೆಂಬಲಿಸು}}, \emph{bembalisu}.) To burst? (\textbf{\kn{ಸಿಡಿ}}, \emph{siḍi}.) To snore? (\textbf{\kn{ಗೊರಕೆ} \kn{ಹೊಡೆ}}, \emph{gorake hoḍe}.) The forehead? (\textbf{\kn{ಹಣೆ}}, \emph{haṇe}.) The throat? (\textbf{\kn{ಗಂಟಲು}}, \emph{gaṇṭalu}.) The heel? (\textbf{\kn{ಹಿಮ್ಮಡಿ}}, \emph{himmaḍi}.) The thigh? (\textbf{\kn{ತೊಡೆ}}, \emph{toḍe}.) The palm? (\textbf{\kn{ಅಂಗೈ}}, \emph{aṅgai}.) Patience? (\textbf{\kn{ತಾಳ್ಮೆ}}, \emph{tāḷme}.) Jealousy? (\textbf{\kn{ಅಸೂಯೆ}}, \emph{asūye}.) Loneliness? (\textbf{\kn{ಒಂಟಿತನ}}, \emph{oṇṭitana}.) Curiosity? (\textbf{\kn{ಕುತೂಹಲ}}, \emph{kutūhala}.) A feeling? (\textbf{\kn{ಭಾವನೆ}}, \emph{bhāvane}.) Childbirth? (\textbf{\kn{ಹೆರಿಗೆ}}, \emph{herige}.) Eyesight? (\textbf{\kn{ದೃಷ್ಟಿ}}, \emph{dṛṣṭi}.) An infection? (\textbf{\kn{ಸೋಂಕು}}, \emph{sōṅku}.) A blister? (\textbf{\kn{ಬೊಕ್ಕೆ}}, \emph{bokke}.) Nutritious? (\textbf{\kn{ಪೌಷ್ಟಿಕ}}, \emph{pauṣṭika}.)
\end{tcolorbox}
\section[(dialogue)]{Words from chapters 171-175 — a second pass}

\label{lesson:KA-R175-second-pass-words-from-chapters-171-175}



\begin{quote}
Say the words for always and now and then.
\end{quote}

\subsection*{Your turn}
Say these aloud:

\begin{itemize}
\raggedright
  \item always, now and then, a date, the weekend, before
  \item rarely, every day, the beginning, a century, a moment
  \item a kitchen, a bathroom, a toilet, a room, a courtyard
  \item a tourist, a lake, a cave, a zoo, a palace
  \item a fence, a pillar, beside, on the other side, on this side
\end{itemize}

\begin{tcolorbox}[breakable,colback=teal!4,colframe=teal!35!black,title={Before you move on}]
Always? (\textbf{\kn{ಯಾವಾಗಲೂ}}, \emph{yāvāgalū}.) Now and then? (\textbf{\kn{ಆಗಾಗ}}, \emph{āgāga}.) A date? (\textbf{\kn{ದಿನಾಂಕ}}, \emph{dināṅka}.) The weekend? (\textbf{\kn{ವಾರಾಂತ್ಯ}}, \emph{vārāntya}.) Before? (\textbf{\kn{ಮುಂಚೆ}}, \emph{muñce}.) Rarely? (\textbf{\kn{ಅಪರೂಪಕ್ಕೆ}}, \emph{aparūpakke}.) Every day? (\textbf{\kn{ಪ್ರತಿದಿನ}}, \emph{pratidina}.) The beginning? (\textbf{\kn{ಆರಂಭ}}, \emph{ārambha}.) A century? (\textbf{\kn{ಶತಮಾನ}}, \emph{śatamāna}.) A moment? (\textbf{\kn{ಕ್ಷಣ}}, \emph{kṣaṇa}.) A kitchen? (\textbf{\kn{ಅಡುಗೆಮನೆ}}, \emph{aḍugemane}.) A bathroom? (\textbf{\kn{ಸ್ನಾನದ} \kn{ಮನೆ}}, \emph{snānada mane}.) A toilet? (\textbf{\kn{ಶೌಚಾಲಯ}}, \emph{śaucālaya}.) A room? (\textbf{\kn{ಕೋಣೆ}}, \emph{kōṇe}.) A courtyard? (\textbf{\kn{ಅಂಗಳ}}, \emph{aṅgaḷa}.) A tourist? (\textbf{\kn{ಪ್ರವಾಸಿ}}, \emph{pravāsi}.) A lake? (\textbf{\kn{ಸರೋವರ}}, \emph{sarōvara}.) A cave? (\textbf{\kn{ಗುಹೆ}}, \emph{guhe}.) A zoo? (\textbf{\kn{ಮೃಗಾಲಯ}}, \emph{mṛgālaya}.) A palace? (\textbf{\kn{ಅರಮನೆ}}, \emph{aramane}.) A fence? (\textbf{\kn{ಬೇಲಿ}}, \emph{bēli}.) A pillar? (\textbf{\kn{ಕಂಬ}}, \emph{kamba}.) Beside? (\textbf{\kn{ಪಕ್ಕ}}, \emph{pakka}.) On the other side? (\textbf{\kn{ಆಚೆ}}, \emph{āce}.) On this side? (\textbf{\kn{ಈಚೆ}}, \emph{īce}.)
\end{tcolorbox}
